\documentclass[10pt]{article}
\usepackage[margin=20mm]{geometry}
\usepackage{amsmath,amssymb,mathtools,fontspec,unicode-math,xurl,hyperref}
\IfFontExistsTF{DejaVu Serif}{\setmainfont{DejaVu Serif}}{\setmainfont{Latin Modern Roman}}
\setmathfont{latinmodern-math.otf}
\setlength{\parindent}{0pt}\setlength{\parskip}{6pt}\sloppy
\title{Quantization in a grain-boundary interpolation:
a technical verification report}
\author{}
\date{}
\begin{document}
\maketitle
Abstract: The acquired 2023 preprint of Fortuna, Garroni and Spadaro is examined with emphasis on the interpolation underlying its grain-boundary cell formula. A constant-input example shows that the displayed construction in Lemma 6.3 does not have the asserted lattice-valued tangential jumps. A recursive residual correction restores the local trace, quantization and energy estimates. This is a demonstrated defect in a construction used by the central proof, rather than a counterexample to the main theorem. A separate local energy estimate omits the common shear component. The findings concern the retained preprint; equivalence with the 2025 journal article has not been established.\par
\section*{Paper identity, version and inspected scope}
The target record is M25-0403, On the Read–Shockley energy for grain boundaries in 2D polycrystals, by Martino Fortuna, Adriana Garroni and Emanuele Spadaro, DOI 10.1002/cpa.22245 (https://doi.org/10.1002/cpa.22245). The actual source is the 43-page institutional copy of arXiv:2306.07742v1, dated 13 June 2023, whose title is On the Read-Shockley energy for grain boundaries in poly-crystals [FGS]. Its numbered pages coincide with PDF pages 1–43. References below to a source equation, theorem or page refer to this version.\par
The source text, including both appendices, was read. Detailed checks concentrated on the circulation constraint, the quadratic-energy compactness argument, the periodic-cell normalization, the compression construction, and the interpolation used to identify the cell energy. The two findings below are supported by complete local calculations. The external logarithmic estimate and the density argument invoked in the source were not independently re-proved. The statement of the incompatible-field rigidity estimate was compared with the retained reference [Retained latex: \textbackslash{}cite[Theorem 3.3, p.\textasciitilde{}8]\{MSZ\}]; its full proof was not independently certified. Available reference PDFs and earlier extraction/rendering records were retained. No author program was executed.\par
\section*{Definitions and the relevant dependency}
The source works with a matrix field \(A\) and a relatively closed defect set \(S\). In addition to \(\operatorname{supp}\operatorname{Curl}A\subset S\), admissibility requires every relevant closed-loop circulation outside \(S\) to lie in \(\tau\varepsilon\mathbb Z^2\), where \(\tau,\varepsilon>0\) [Retained latex: \textbackslash{}cite[p.\textasciitilde{}4, (H2)]\{FGS\}]. The energy is the sum of a quadratic distance from \(\mathrm{SO}(2)\) outside the core and the core area, divided by \(\varepsilon\).\par
Equation reference: eq:M. For the interpolation lemma, put \[Q_{\ell,L}^0=(0,\ell)\times(0,L),\qquad
 G=\ell\int_0^L |g(t)|^2\,dt,\qquad
 q=\frac{\tau\varepsilon}{\ell},\] where \(L,\ell>0\), \(\ell\ge\varepsilon\) and \(g\in L^2_{\rm loc}(\mathbb R;\mathbb R^2)\). The source additionally assumes \(G<\varepsilon L\) for its quantitative estimate and sets \[
 M=\left(\frac{L\varepsilon}{G}\right)^{1/3},\qquad
 N=\left\lfloor\frac{L}{M\ell}\right\rfloor+1.\] The case \(g=0\) is treated separately. Each complete strip \(E_i\) has height \(M\ell\), and the final strip has the remaining height. On a strip with bottom coordinate \(y_i\), the local primitive is \[u_i(x_1,x_2)=
 \left(1-\frac{x_1}{\ell}\right)\int_{y_i}^{x_2}g(t)\,dt,\qquad A_i=\nabla u_i.\] Thus \(A_i e_1=0\) on its bottom and \[A_i e_1=\bar g_i:=-\ell^{-1}
       \int_{y_i}^{y_i+M\ell}g(t)\,dt\] on its top when the strip is complete.\par
Equation reference: eq:printed. The source defines \(g_i^\varepsilon=P_{\tau\varepsilon\mathbb Z^2}(\ell\bar g_i)/\ell\), using componentwise floor projection, and assembles \[
 A=A_1\chi_{E_1}
   +\sum_{i=2}^{N}\chi_{E_i}
       \bigl[A_i+(g_{i-1}^{\varepsilon}-\bar g_{i-1})\otimes e_1\bigr].\] The assertion immediately before source (6.13) is that the jumps of \(Ae_1\) belong to \(q\mathbb Z^2\). This is the property examined next.\par
\section*{A failed quantization step and a local correction}
Problem 1 (The offsets in Lemma 6.3 do not preserve the lattice jumps).  Equation [eq:printed] (\#eq:printed), printed on source p. 33 immediately before (6.13), does not generally have the asserted \(q\mathbb Z^2\)-valued jumps.\par
\subsection*{An example satisfying all hypotheses}
Equation reference: eq:parameters, eq:failure. Choose \[
 \tau=1,\quad \varepsilon=\ell=\frac1{10},\quad L=1,\quad
 g(t)=\frac18e_1\quad (t\in\mathbb R).\] The function is locally square integrable and nonzero, \(\ell=\varepsilon\), and \[G=\frac1{640}<\frac1{10}=\varepsilon L.\] Equations [eq:M] (\#eq:M) give \(M=4\), \(N=3\). The strip heights are \(2/5,2/5,1/5\), and \(q=1\). For the first two strips, \[\bar g_1=\bar g_2=-\frac12e_1,\qquad
 g_1^\varepsilon=g_2^\varepsilon=-e_1.\] Consequently the offsets in [eq:printed] (\#eq:printed) are \[(c_1,c_2,c_3)=\left(0,-\frac12e_1,-\frac12e_1\right).\] At the second interface, \(x_2=4/5\), the lower trace is \(\bar g_2+c_2=-e_1\), whereas the upper trace is \(c_3=-e_1/2\). Therefore \[
 [Ae_1]=\frac12e_1\notin\mathbb Z^2=q\mathbb Z^2.\] The first interface happens to have zero jump; this does not repair the second one. The calculations are exact rational arithmetic. An independently written symbolic check reproduces them; its source and actual output are retained with the verification files.\par
The failure persists at arbitrarily small scales: set \(\ell=\varepsilon\), \(L=1\), retain \(g=e_1/8\), and take \(0<\varepsilon<1/8\). Then \(M=4\), the first two complete strips exist, and their second-interface jump is still \(e_1/2\), whereas \(q=1\).\par
\subsection*{Alternative readings}
Reversing the orientation of the interface changes the sign in [eq:failure] (\#eq:failure), without making the jump integral. The source explicitly sets the lower trace of each local primitive to zero, so interpreting \(u_i\) as a single unshifted global primitive contradicts its displayed endpoint conditions. Its projection is expressly componentwise floor projection. Changing a rounding choice by lattice vectors cannot remove the half-integer residue.\par
There are two connected defects in [eq:printed] (\#eq:printed): the residual is inserted with the wrong sign for the first matching operation, and later interfaces require the previous residual to be carried into the next projection. A sign change alone is not a general repair.\par
\subsection*{A corrected interpolation}
Equation reference: eq:recurrence. Let \(P_q\) be componentwise floor projection onto \(q\mathbb Z^2\). Set \[
 c_1=0,\qquad
 c_{i+1}=c_i+\bar g_i-P_q(c_i+\bar g_i),\qquad
 \widetilde A|_{E_i}=A_i+c_i\otimes e_1.\] For the example [eq:parameters] (\#eq:parameters), this gives \((c_1,c_2,c_3)=(0,e_1/2,0)\), and the two jumps are \(e_1\) and \(0\).\par
Proposition 1.  Under the hypotheses of Lemma 6.3, the replacement [eq:recurrence] (\#eq:recurrence) preserves its prescribed second column, bottom and top first-column conditions, lattice-valued jumps, and the energy order in source (6.11). Constants may depend on the fixed parameters \(\lambda,\tau\).\par
Proof. At interface \(i\), the jump is \[c_{i+1}-(c_i+\bar g_i)
   =-P_q(c_i+\bar g_i)\in q\mathbb Z^2.\] Every \(c_i\) belongs to \([0,q)^2\), so \(|c_i|<\sqrt2q\). Adding the constant matrix \(c_i\otimes e_1\) changes neither the second column nor the interior curl of \(A_i\). The bottom first column remains zero because \(c_1=0\); the top first column is constant because the primitive’s integral and \(c_N\) are constant there. The only possible new curl lies on the finitely many horizontal strip interfaces.\par
Equation reference: eq:elastic. For a strip of height \(h_i\le M\ell\), the elementary estimate corresponding to source (6.8) yields \[\int_{E_i}|A_i|^2
 \le \left(\frac{h_i^2}{\ell^2}+\frac13\right)
        \ell\int_{y_i}^{y_i+h_i}|g|^2.\] Hence \[
 \int_{Q_{\ell,L}^0}|\widetilde A|^2
 \le 2\left(M^2+\frac13\right)G
       +4\tau^2\frac{\varepsilon^2L}{\ell}.\] There are at most \(N-1\le L/(M\ell)\) interfaces, each of length \(\ell\). Their tubular neighborhoods therefore satisfy \[|B_{\lambda\varepsilon}(S)|
 \le \frac{L}{M\ell}
       \bigl(2\lambda\varepsilon\ell+\pi\lambda^2\varepsilon^2\bigr)
 \le (2\lambda+\pi\lambda^2)\frac{\varepsilon L}{M},\] where \(\varepsilon\le\ell\) was used. Overlaps or truncation by the ambient rectangle only reduce this upper bound. Since \(G<\varepsilon L\), \(M>1\), and \[M^2G=\frac{\varepsilon L}{M}
       =\varepsilon^{2/3}L^{2/3}G^{1/3},\] combining this identity with [eq:elastic] (\#eq:elastic) proves \[\int_{Q_{\ell,L}^0}|\widetilde A|^2+|B_{\lambda\varepsilon}(S)|
 \le C_{\lambda,\tau}
       \varepsilon^{2/3}L^{2/3}G^{1/3}
       +C_{\lambda,\tau}\frac{\varepsilon^2L}{\ell}.\] This is the claimed local order. Empty terminal strips, if \(L/(M\ell)\) is an integer, may simply be omitted. ◻\par
\subsection*{Dependency and impact}
The source invokes Lemma 6.3 in the proof of Proposition 6.2 and uses it to assert admissibility after (6.21), on p. 35. Section 6.3 then explicitly applies Proposition 6.2 to prove Theorem 5.3. That theorem supplies the identification in Corollary 5.4 and the cell formula used in Section 7.\par
The demonstrated failure therefore lies on a central proof dependency and is assessed as Major under the stated core-contribution criterion. The assessment concerns the invalid displayed construction. Proposition Proposition 1 (\#prop:repair) supplies a complete local replacement, which is counterevidence against treating the example as a false-theorem result. The complete downstream proof has not been independently reconstructed after this replacement, and no statement about an author-adopted correction or the journal-final version is made.\par
\section*{A local energy bound loses the common shear}
Problem 2 (The bound in source (5.21)–(5.22) omits a mean component).  For the general compression-to-compression parameters in Section 5.2.1, the right-hand side of source (5.21) can vanish although its left-hand side is strictly positive.\par
Choose \(\ell=h=1\), \(r=\rho=1/4\), \(\eta=\widetilde\eta=1\) and \(\mu=\widetilde\mu=1/2\). These satisfy the rectangular size restrictions on source p. 20 and keep the denominators in the construction nonzero. Source (5.15) reduces, on both interpolation regions, to \[\widehat D_i=
 \begin{pmatrix}1&1/2\\0&1\end{pmatrix},\qquad i=3,4.\] All difference terms on the right-hand side of (5.21) vanish, and \(\eta+\widetilde\eta-2=0\). Its claimed upper bound is therefore zero. However, \[|\widehat D_i-I|^2=\frac14,\qquad
 \operatorname{dist}^2(\widehat D_i,\mathrm{SO}(2))=\frac{17}{4}-\sqrt{17}>0.\] For completeness, minimizing the squared Frobenius distance over rotations gives \(\operatorname{dist}^2(B,\mathrm{SO}(2))=|B|^2+2-2\sqrt{(b_{11}+b_{22})^2+(b_{21}-b_{12})^2}\); substitution gives the displayed value. Thus this is a failure of the actual inequality, not a change of norm or orientation. The corresponding integral bound (5.22) also fails because the interpolation regions have positive area. In this example the circulation around the core is zero by source (5.17), so quantization imposes no excluded parameter condition.\par
Equation reference: eq:mean-bound. A valid replacement follows directly from the printed gradients. Define \[d=(\widetilde\mu-\mu,\widetilde\eta-\eta),\qquad
 a=\left(\frac{\mu+\widetilde\mu}{2},
          \frac{\eta+\widetilde\eta}{2}-1\right).\] For \(\Delta_3\), the first column of \(\widehat D_3-I\) is \(\frac{h}{2\ell}(1-h/x_2)d\), while the second is \(a+\frac{h^2x_1}{2\ell x_2^2}d\). The corresponding sign in the first column changes on \(\Delta_4\). Applying \(|v+w|^2\le2|v|^2+2|w|^2\) gives \[
 \operatorname{dist}^2(\widehat D_i,\mathrm{SO}(2))
 \le C|d|^2\left[
       \frac{h^2}{\ell^2}
       +\frac{h^4}{\ell^2}
         \left(\frac1{x_2^2}+\frac{x_1^2}{x_2^4}\right)\right]
       +2|a|^2.\] This retains the missing mean shear and a nonnegative squared mean compression. Merely squaring the source’s last term would still omit the common shear.\par
The specific parameters in source (5.37) satisfy \(\widetilde\mu=-\mu\) and \(\eta+\widetilde\eta=2\cos\alpha\). In that specialization, \(|a|^2=(1-\cos\alpha)^2=O(\alpha^4)\). Summing its integrated contribution over the compression cells, using \(L>h_D\) and \(\ell_D=\varepsilon\tau/\sin\alpha\), gives an additional \(O(L\varepsilon\tau\alpha^3)\) for positive small \(\alpha\). This is below the target \(L\varepsilon\tau\alpha(|\log\alpha|+1)\) order. Thus the counterexample establishes a Minor defect in the general local estimate; it does not demonstrate failure of the specialized Read–Shockley upper bound.\par
\section*{Outcome and remaining scope}
The highest supported impact in the acquired source is Major, from Problem 1’s central construction. Problem 2 has a narrower, local impact. The verified repairs concern the particular interpolation and pointwise estimate, respectively. Neither establishes author adoption or substitutes for a check of a later source version.\par
The literal extension in Remark 2.8 to an energy comparable to \(\operatorname{dist}(A,\mathrm{SO}(2))\), rather than its square, was identified but is not settled by the quadratic-energy analysis here. A difference between the energies along one sequence would not itself refute a \(\Gamma\)-limit. Likewise, the full treatment of periods around holes in the general domain is outside the repaired rectangular interpolation result. These unresolved scope points do not increase the supported severity.\par
Item | Assessment | Repair verified here | Adoption\par
Lemma 6.3 construction | Major | Local traces, lattice jumps and energy order | Unknown\par
Source (5.21)–(5.22) | Minor | Pointwise bound with mean terms; specialized extra order | Unknown\par
Scientific decision | Pending | Independent acceptance remains separate | No human review\par
\section*{Supporting checks}
The main source PDF has SHA256 \[\text{\texttt{3dcfcd72fcf62926793181f23f3ca9dd3ac8d}}\] \[\text{\texttt{4efdef11b1ec387e5ccf79f69fc}}.\] The split above is typographical only; the verification record stores the uninterrupted 64-character value. Calculations in the report are analytic arguments; the finite symbolic script checks the displayed example and algebra, not the entire main theorem.\par
The report was produced in this recovery session. The controller supplied a Personal/Chat/GPT-6/Pro interface observation; backend equivalence was not independently established. No human-review declaration is supplied.\par
9 M. Fortuna, A. Garroni and E. Spadaro, On the Read-Shockley energy for grain boundaries in poly-crystals, arXiv:2306.07742v1, 13 June 2023, 43-page Sapienza institutional copy. Target journal record: DOI 10.1002/cpa.22245 (https://doi.org/10.1002/cpa.22245). S. Müller, L. Scardia and C. I. Zeppieri, Geometric rigidity for incompatible fields and an application to strain-gradient plasticity, retained 26-page preprint, Theorem 3.3 on p. 8. Published bibliographic record cited in the acquired main source: Indiana Univ. Math. J. 63 (2014), 1365–1396.\par
\end{document}
